\subsection{Numerical Scope and Scientific Limitations of the Eulerian Families}
\label{sec:eulerian-family-scope}

The CPU implementation offers three complementary Eulerian descriptions of
spherical kinetic evolution.  Each stores and transports the distribution in
a different form and treats phase-space boundaries accordingly.  The choice
therefore depends on the orbits, boundary fluxes, and collision processes that
the calculation must resolve.

Here \emph{Eulerian} means that the distribution is stored on a fixed
phase-space mesh, rather than carried by individual moving markers.  An
\emph{upwind} face flux takes its distribution value from the side of a cell
face from which the characteristic flow arrives; the first-order rule uses
the nearest stored value on that upstream side.

The canonical family uses \(\eta=L^2/2\), the angular-momentum coordinate;
\(H=\Phi(r)+w^2/2+\eta/r^2\), the one-particle Hamiltonian in a frozen
potential \(\Phi\); and \(\tau_{\rm orb}\in[0,1)\), the radial-orbit phase
normalized by one radial period.  The subscript distinguishes this orbit
phase from the Schwarzschild time \(\tau\) used in the relativistic chapters.
The canonical chapter denotes this normalized orbit phase by
\(\zeta_{\rm orb}\).  In a relativistic context, \(\tau\) denotes
Schwarzschild time.  The orbit atlas tabulates the represented bound orbit families.
An \emph{open packet} is a continuous phase-space carrier for matter that the
bound atlas cannot currently represent, such as an unbound or otherwise
unsupported orbit; VENUS retains and advances that matter instead of silently
discarding it.

\begin{center}
\footnotesize
\begin{tabular}{@{}>{\raggedright\arraybackslash}p{0.19\linewidth}
                  >{\raggedright\arraybackslash}p{0.38\linewidth}
                  >{\raggedright\arraybackslash}p{0.34\linewidth}@{}}
\toprule
Solver family & Numerical representation & Scientific limitations and suitable use \\
\midrule
Explicit finite volume &
Cell averages on the spherical \((r,w,L)\) or \((r,w,\eta)\) mesh, advanced
with first-order upwind radial and radial-velocity face fluxes.  The
implementation includes reflecting, vacuum-outflow, and prescribed-inflow
faces, threaded flux construction, and a multidimensional CFL limit. &
The first-order stencil is diffusive near sharp orbit boundaries.  Long-time
equilibria require mesh refinement in every occupied phase-space direction,
and boundary fluxes must be included explicitly in mass and energy budgets. \\
\addlinespace
Grid semi-Lagrangian &
Nodal characteristic pullback for EV and VP, plus cell-integrated
conservative and Hamiltonian SLDG for VP.  The conservative path uses
positivity-limited reconstruction, mapped-cell overlap integrals, midpoint
self-gravity, recursive map refinement, and a delayed central carrier. &
The nodal path is not locally conservative.  The cell-integrated Hamiltonian
path is restricted to a static background, and very low angular momentum,
long delayed returns, rapidly changing fields, and map fallbacks require
problem-specific refinement. \\
\addlinespace
Canonical orbit phase &
A bound-orbit atlas in \((\eta,H,\tau_{\rm orb})\), frozen-potential phase translation,
Eulerian overlap maps, field rebuilds, finite-angle collision couplings, open
packets, positive packet compression, independent event clocks, and
versioned restarts. &
The representation is most efficient when the potential changes slowly
compared with an orbital period.  Unbound tails remain explicit packets, and
the orbit atlas, continuous-packet population, collision locality, and field
refresh cadence must all be refined together. \\
\bottomrule
\end{tabular}
\end{center}

For the finite-volume update, a cell \((i,j,k)\) loses the difference of its
two radial-face fluxes and its two radial-velocity-face fluxes.  This is a
conservative balance law: a reflecting face returns the appropriate
velocity-space flux, while an outflow or prescribed-inflow face contributes
to the recorded boundary budget.  Adding the interior matter to those
boundary transfers gives the conserved total.  This complete budget is the
quantity to check after each step.

The conservative SLDG family replaces face fluxes by a forward map of each
source cell.  It limits a local polynomial to remain nonnegative, subdivides
a mapped patch when its area or reversibility defect is too large, integrates
exact polynomial moments over polygon--cell intersections, and records mass
that cannot be represented by the accepted map.  Midpoint gravity and the
persistent delayed-inner carrier keep field evolution and unresolved central
passages in the same conservative accounting.  These ingredients are
described in detail in Section~\ref{sec:cell-integrated-sldg}.

The canonical family changes coordinates instead.  For a frozen spherical
potential, angular momentum \(\eta\), Hamiltonian \(H\), and orbital phase
\(\tau_{\rm orb}\) separate stationary orbit structure from phase translation.  Matter
created by a collision is retained as a continuous packet until it can be
mapped positively onto the orbit atlas; bound and open-orbit matter are never
silently conflated.  The full construction is described in
Section~\ref{sec:fully-kinetic-bp}.

All three families require a numerical convergence study appropriate to the
chosen physical problem.  Refine every coordinate or sampling budget that
controls the active operator; compare density, force, characteristic
invariants, and collision moments; reduce the timestep and any operator
splitting interval; verify that boundary and hidden-carrier ledgers close;
and compare uninterrupted and restarted evolution.  Supported OpenMP and MPI
layouts should reproduce the same physical observables within the declared
floating-point and stochastic uncertainty.  For collisional calculations,
also require collision regions smaller than the local mean free path and
separate Monte--Carlo scatter from deterministic discretization error.
